\documentclass[10pt,a4paper]{amsart}
\usepackage[margin=19mm]{geometry}
\usepackage{amsmath,amssymb,mathtools}
\usepackage[scr=rsfs,cal=euler]{mathalfa}
\usepackage{xfrac}
\usepackage[unicode,hidelinks,bookmarks=false]{hyperref}
\newcommand{\ie}{i.e.\ }
\newcommand{\cf}{cf.\ }
\newcommand{\resp}{respectively\ }
\newcommand{\Subsection}{\subsection}
\providecommand{\nobreakdash}{\nobreak\hbox{-}}
\setlength{\emergencystretch}{2em}
\multlinegap0pt
\usepackage{tikz}
\usetikzlibrary{cd}
\usepackage{mathtools}
\usepackage{amssymb}
\usepackage{xcolor}
\newtheorem{lemma}{Lemma}
\newtheorem{theorem}{Theorem}
\def\Ext{\text{Ext}}
\def\Z{{\mathbb{Z}}}
\title{Arbitrary models of the complete first-order theories of FDZ-rings}
\author{Mahmood Sohrabi}
\date{Source: arXiv:2603.27730v1 (CC BY 4.0)}
\begin{document}
\maketitle

\noindent\emph{Source and licence.} Arbitrary models of the complete first-order theories of FDZ-rings. Version arXiv:2603.27730v1 (CC BY 4.0), distributed by arXiv under the Creative Commons Attribution 4.0 International licence, http://creativecommons.org/licenses/by/4.0/, as recorded in the arXiv licence metadata for this exact version and shown on the abstract page https://arxiv.org/abs/2603.27730v1. Full licence notice: \texttt{licenses/arxiv-2603.27730-CC-BY-4.0.txt}.\par\smallskip

\noindent\textit{Editorial scope.} Counted unit: the article's theorem theorem (source0.tex lines 661--664). The article's complete original proof of it is delivered with it (source0.tex lines 665--719, one \texttt{proof} environment of 55 lines). No mathematics is written, repaired or re-derived: the statement and the proof are the article's own lines, in the article's own notation, and no retained line is substituted or re-flowed.  Retained uncounted as the source-local context the proof needs: lines 540--568; lines 585--585.  Nothing was omitted from any retained span, so the package carries no omission record.  No retained line contains a citation, so the wrapper carries no bibliography and no bibliography entry.  The retained lines contain the article's own diagram code, which the wrapper typesets with the standard tikz packages; no image file is delivered and the package declares no asset dependency.  Editorial formatting only, in addition: an AMS \texttt{amsart} preamble that restates the article's own declarations and macros used by the retained lines, namely the environments \texttt{lemma}, \texttt{theorem} on separate counters, together with the standard packages those lines need.  The retained environments print in this document as Lemma 1, Theorem 1 in order of appearance, whereas the article prints the same items with the numbering of its own full document.  The complete native source of the article as distributed in the arXiv e-print for this exact version is bundled unchanged as \texttt{sources/197381/source0.tex}.
\begin{lemma}\label{lem:6term}
    Let $A$ be an FDZ-ring where $\text{Spec}^0(P(f_A))$ is connected, and let $B$ be a ring that is elementarily equivalent to $A$. Then there exists a scalar ring $R\equiv \Z$ such that the following is true:
    \begin{enumerate}
 \item  $Ann(B)$ fits into the following short exact sequence of abelian groups
    $$ 0\to O(A(R)) \to Ann(B)\to B_0\to 0$$
    where $O(A(R))=Ann(A(R))\cap \Delta(A(R))$ and $B_0\equiv A_0$ are abelian groups.
\item There exist isomorphisms, $\eta, \psi, \phi$ and $\mu$ of $R$-algebras: %ring isomorphisms $\psi: \widehat{A(R)}\to \widehat{B}$ and $\phi: \Delta(A(R))\to \Delta(B) $ 
that make the following diagram commutative:

   
$$\begin{tikzcd}[column sep=small]
0 \arrow[r] & O(A(R)) \arrow[r] \arrow[d,"\eta"] &
\Delta(A(R)) \arrow[r] \arrow[d,"\psi"] &
\widehat{A(R)} \arrow[r] \arrow[d,"\phi"] &
A/K(A(R)) \arrow[r] \arrow[d,"\mu"] & 0 \\
0 \arrow[r] & O(B) \arrow[r] &
\Delta(B) \arrow[r] &
\widehat{B} \arrow[r] &
B/K(B) \arrow[r] & 0
\end{tikzcd}$$

where the horizontal arrows represent the restrictions of the obvious canonical mappings.
   \item $B$ fits into the following short exact sequence of rings
    $$0\to Ann(B) \to B \to \widehat{A(R)}\to 0$$ 

   
    \end{enumerate}
     
\end{lemma}

\subsection{Construction of abelian deformations using symmetric 2-cocycles} \label{subsec:cocycles}

\begin{theorem}[Characterization Theorem] Let $A$ be an FDZ-ring where $\text{Spec}^0(P(A))$ is connected. If $B$ is any ring $B\equiv A$, then
$B\cong (A(R),B_0, f,g,h)$
for some ring $R\equiv \Z$, some $B_0\equiv A_0$, and choice of cocycles $f,g$ and $h$ as described in Section~\ref{subsec:cocycles}.   
\end{theorem}

\begin{proof}
    We have already proven that $B\equiv A$ satisfies conditions (1.)-(3.) of  Lemma~\ref{lem:6term}. We shall refer to the maps and constructions mentioned in that lemma. Let us denote $(A(R),B_0,f,g,h)$ by $C$, where $f$, $g$ and $h$ are yet to be defined in terms of those defining the structure of the ring $B$. By construction, there is a cocycle $h'$ that defines the lower short exact sequence of the following diagram. Therefore, there exists a cocycle $h\in S^2(B_0, O(A(R))$ defining the upper exact sequence which induces an isomorphism $Ann(C)\to Ann(B)$ making the following diagram commutative.
$$\begin{tikzcd}[column sep=small]
0 \arrow[r] & O(A(R)) \arrow[r]  \arrow[d,"\eta"] &
Ann(C) \arrow[r] \arrow[d]&
B_0 \arrow[r]\arrow[d,"id"] &0\\
0 \arrow[r] & O(B) \arrow[r] &
Ann(B) \arrow[r]  &
B_0 \arrow[r]  & 0
\end{tikzcd}$$
Since $\Delta(C)=\Delta(A(R))$ and the isomorphism $\psi: \Delta(A(R))\to \Delta(B)$ agrees with the isomorphism $\eta$ and $K(B)=\Delta(B)+Ann(B)$, in fact $K(B)$ is defined as an extension of $\Delta(B)$ by $B_0$ using the cocycle $h'$, therefore, the cocycle $h$ defines $K(C)$ in a manner that there exists an isomorphism $\theta:K(C) \to K(B)$ making the following diagram commutative:

$$\begin{tikzcd}[column sep=small]
0 \arrow[r] & \Delta(A(R)) \arrow[r]  \arrow[d,"\psi"] &
K(C) \arrow[r] \arrow[d, "\theta"]&
B_0 \arrow[r]\arrow[d,"id"] &0\\
0 \arrow[r] & \Delta(B) \arrow[r] &
K(B) \arrow[r]  &
B_0 \arrow[r]  & 0
\end{tikzcd}$$

We note that $C/K(C)$ can be identified with $A(R)/K(A(R))$ and the isomorphism $\mu:C/K(C) \to B/K(B)$ respects the maximal torsion subgroup $N$. Therefore, $B/K(B)\cong \mu(N) \oplus M(B)$, where $M(B)\cong M(A(R))=M(C)$ as free $R$-modules. If one picks an $R$-basis, $m_1, \ldots ,m_k$ for $M(C)$ and the corresponding basis $\mu(m_1), \ldots \mu(m_k)$, then the inverse images of these under the corresponding canonical epimorphisms $\pi:C\to\widehat{C}$ and $\pi':B\to \widehat{B}$ generate subgroups $M_1(C)$ and $M_1(B)$ of $C$ and $B$, respectively, which contain the corresponding annihilators. Therefore, given a cocycle $f'\in S^2(M(B), Ann(B)$, defining the extension of in the lower row of the following diagram, there exists a cocycle $f\in S^2(M(C), Ann(C)$ defining the extension in the upper row, and an isomorphism $\theta_1$ making the diagram commutative:
$$\begin{tikzcd}[column sep=small]
0 \arrow[r] & K(C) \arrow[r]  \arrow[d,"\theta"] &
K(C)+M_1(C) \arrow[r] \arrow[d, "\theta_1"]&
M(C) \arrow[r]\arrow[d,"\mu|_{M(C)}"] &0\\
0 \arrow[r] & K(B) \arrow[r] &
K(B)+M_1(B) \arrow[r]  &
M(B) \arrow[r]  & 0
\end{tikzcd}$$
Since we have 
\begin{equation}\label{eqn:exts}
\begin{split}
\Ext(B/K(B),K(B))&\cong \Ext(M(B)\oplus N(B), K(B))\\
&\cong  \Ext(M(B), K(B))\oplus  \Ext(N(B), K(B))
\end{split}
\end{equation}

it is left to consider the extension:
$$0\to K(B)\to L(B) \to N(B)\to 0$$
Since $N(B)$ is the maximal torsion subgroup of $B/K(B)$, the cocycle, say $g'$, whose class in $\Ext (N(B), K(B))$ partially determines $L(B)$ will produce the subgroup $L(B)=Is(Ann(B) + B^2)$. Consider the corresponding extension in $A(R)$:
$$0\to K(A(R))\to L(A(R)) \to N\to 0$$
and let $q\in S^2(N, K(A(R))$ be the cocycle defining this extension up to equivalence and
Let $\pi:A(R) \to \widehat{A(R)}$ be the canonical epimorphism and $\pi': B \to \widehat{B}$ the one for $B$. The homomorphism $\pi$ induces a homomorphism: $$\pi_*: \Ext(N, K(A(R)) \to \Ext(N, \pi(K(A(R))).$$ Identifying $\widehat{A(R)}$ and $\widehat{B}$ via the isomorphism $\phi$, there exists a cocycle $p'\in \Ext(N(B), K(B))$ such that $\pi'_*(p')=\pi_*(q)$. It is also clear that for any choice of $g'\in S^2(N(B), Ann(B))$, we have $p'+g' \in S^2(N, K(B))$ and $\pi_*'(p'+g')=\pi_*'(p')=\pi_*(q)$. Finally, it is clear that there exists a cocycle $p\in S^2(N, K(C))$, and some cocycle $g\in S^2(N, Ann(C))$ such that $\pi''_*(p+c)=\pi'_*(p'+g')$, where $\pi'':C\to \widehat{C}$ is the canonical epimorphism and an isomorphism $\theta_2: L(C) \to L(B)$, making the following diagram commutative
$$\begin{tikzcd}[column sep=small]
0 \arrow[r] & K(C) \arrow[r]  \arrow[d,"\theta"] &
L(C) \arrow[r] \arrow[d, "\theta_2"]&
N \arrow[r]\arrow[d,"\mu|_{N}"] &0\\
0 \arrow[r] & K(B) \arrow[r] &
L(B) \arrow[r]  &
N(B) \arrow[r]  & 0
\end{tikzcd}$$

This finishes the proof considering the isomorphism in Equation~\ref{eqn:exts} above.
\end{proof}

\end{document}
